\documentclass[12pt, a4paper]{article}

% ─── Packages ────────────────────────────────────────────────────────────────
\usepackage[T1]{fontenc}
\usepackage[utf8]{inputenc}
\usepackage[english]{babel}
\usepackage{geometry}
\usepackage{graphicx}
\usepackage{amsmath}
\usepackage{amssymb}
\usepackage{siunitx}
\usepackage{booktabs}
\usepackage{caption}
\usepackage{subcaption}
\usepackage{float}
\usepackage{hyperref}
\usepackage{fancyhdr}
\usepackage{parskip}
\usepackage{xcolor}
\usepackage{array}
\usepackage{multirow}
\usepackage{enumitem}
\usepackage{tcolorbox}
\usepackage{tabularx}
\usepackage{microtype}

\geometry{top=2.5cm, bottom=2.5cm, left=2.5cm, right=2.5cm}

\sisetup{per-mode=symbol, detect-weight=true}

% ─── Placeholder commands ────────────────────────────────────────────────────
\newcommand{\val}[1]{\textcolor{blue}{\textbf{??\,\si{#1}}}}
\newcommand{\valx}{\textcolor{blue}{\textbf{??}}}
\newcommand{\enter}[1]{\textcolor{blue}{\textit{[#1]}}}

\tcbuselibrary{skins}
\newcommand{\hint}[1]{%
  \begin{tcolorbox}[
    enhanced,
    colback  = gray!7,
    colframe = gray!35,
    fontupper= \small,
    left=6pt, right=6pt, top=4pt, bottom=4pt,
    title    = \textbf{\small Hints for this section},
    attach boxed title to top left={yshift=-2mm, xshift=4mm},
    boxed title style={colback=gray!35, colframe=gray!35,
                       fontupper=\small\bfseries, left=3pt, right=3pt}
  ]
  #1
  \end{tcolorbox}}

\hypersetup{colorlinks=true, linkcolor=black, citecolor=black, urlcolor=blue,
            pdftitle={Lab Report - Sequential Circuits},
            pdfauthor={Rupp Arian, Lampert Mathias}}

\pagestyle{fancy}
\fancyhf{}
\lhead{\small EIT \& MEC -- Digi Lab 2: Sequential Circuits}
\rhead{\small Fachhochschule Vorarlberg}
\cfoot{\small Page \thepage}
\lfoot{\small Rev.\,1.0 \quad 05/06.05.2026}
\renewcommand{\headrulewidth}{0.4pt}
\renewcommand{\footrulewidth}{0.4pt}

% ═════════════════════════════════════════════════════════════════════════════
\begin{document}
% ═════════════════════════════════════════════════════════════════════════════

% ─── Title page ──────────────────────────────────────────────────────────────
\begin{titlepage}
  \centering
  \vspace*{3cm}
  {\LARGE\bfseries Laboratory Report\par}
  \vspace{0.5cm}
  {\large\bfseries EIT \& MEC -- Digi Lab 2: Sequential Circuits\par}
  \vspace{0.3cm}
  {\large Fachhochschule Vorarlberg GmbH\par}
  \vspace{2cm}
  \begin{tabular}{ll}
    \textbf{Students:}          & Rupp Arian, Lampert Mathias \\[4pt]
    \textbf{Study programme:}   & Elektronik und Informationstechnologie Dual \\[4pt]
    \textbf{Date of experiment:}& 05/06.05.2026 \\[4pt]
    \textbf{Supervisor:}        & Horatiu O.\,Pilsan \\[4pt]
    \textbf{Report guide:}      & C.\,Thorrold, \textit{Guide for Writing Lab Reports}, v1.0 \\
  \end{tabular}
  \vfill
  {\small Submitted: \today}
\end{titlepage}

% ─── Abstract ────────────────────────────────────────────────────────────────
\begin{abstract}
This report documents the experimental investigation of sequential circuits, conducted
in Digi Lab~2 at Fachhochschule Vorarlberg on 05/06~May~2026.
Three circuits were designed, simulated in LTSpice, and assembled on a breadboard:
a chaser (moving) light based on an 8-bit shift register (74HC164), an extended
version with a ``catch'' function using a D-type flip-flop (74HC74) and a monoflop
(74HC221), and a further extension adding a BCD decade counter (74HC160) with
4-LED display.
The correct behaviour of each circuit was verified by observation on the breadboard,
and the design decisions are discussed and explained.
\enter{Refine after experiments are completed -- max.\ 200 words.}
\end{abstract}

\newpage
\tableofcontents
\newpage

% ─────────────────────────────────────────────────────────────────────────────
\section{Introduction}
% ─────────────────────────────────────────────────────────────────────────────

Sequential circuits, unlike purely combinational circuits, have memory: their output
depends not only on the current inputs but also on the history of previous inputs.
They are the basis of all digital systems that store state, count events, or
synchronise operations to a clock signal.

This laboratory exercise investigates three fundamental sequential building blocks:
the \textbf{shift register}, the \textbf{monoflop}, and the \textbf{counter}.
A shift register stores a sequence of bits and moves them one position on each clock
edge.
A monoflop (one-shot) generates a pulse of defined duration in response to a trigger.
A synchronous BCD counter increments its binary output by one on each rising clock
edge, resetting to zero after reaching nine.

All three components are combined in a progressive design task: first a simple
chaser light is implemented, then extended with a catch function, and finally
augmented with a catch counter displayed in BCD code.

The objectives of this laboratory are:
\begin{itemize}[leftmargin=1.5em]
  \item To understand the operating principle of shift registers, monoflops, and
        synchronous counters.
  \item To design, simulate (LTSpice), and assemble three sequential circuits on
        a breadboard.
  \item To implement a power-on reset circuit ensuring a defined initial state.
  \item To debounce a push-button using a D-type flip-flop.
  \item To combine multiple ICs into a functional system and verify correct
        end-to-end behaviour.
\end{itemize}

The theoretical background was reviewed according to~\cite{Floyd2006}
and~\cite{TietzeSchenk1999}.

% ─────────────────────────────────────────────────────────────────────────────
\section{Equipment and Materials}
% ─────────────────────────────────────────────────────────────────────────────

All circuits were assembled on a standard experimentation (breadboard) system.
A square-wave signal of approximately \SI{1}{\hertz}, \SI{5}{\volt_{pp}}, and
\SI{2.5}{\volt} DC offset from the waveform generator was used as clock input
for all tasks.
A Schmitt-trigger inverter (74HC132) was inserted between the waveform generator
and the logic circuit to ensure correct signal levels.
Table~\ref{tab:equipment} lists the equipment and ICs used.

\begin{table}[H]
  \centering
  \caption{Equipment and components used in this experiment.}
  \label{tab:equipment}
  \begin{tabular}{llc}
    \toprule
    Item & Specification / Part No. & Asset / S/N \\
    \midrule
    Digital oscilloscope          & Teledyne LeCroy HD 4096           & --- \\
    Oscilloscope probe (10:1)     & ---                               & --- \\
    Waveform generator            & \enter{Model}                     & --- \\
    DC power supply               & \SI{5}{\volt}                     & --- \\
    Experimentation board         & ---                               & --- \\
    74HC00 (4x 2-in NAND)         & 74HC00                            & --- \\
    74HC74 (2x D-type flip-flop)  & 74HC74                            & --- \\
    74HC132 (4x NAND Schmitt)     & 74HC132                           & --- \\
    74HC160 (4-bit decade counter)& 74HC160                           & --- \\
    74HC164 (8-bit shift register)& 74HC164                           & --- \\
    74HC221 (2x monoflop)         & 74HC221                           & --- \\
    LEDs (8x for chaser)          & ---                               & --- \\
    LEDs (4x for BCD display)     & ---                               & --- \\
    Push-button                   & ---                               & --- \\
    Resistors, capacitors         & see individual tasks              & --- \\
    \bottomrule
  \end{tabular}
\end{table}

% ─────────────────────────────────────────────────────────────────────────────
\section{Theoretical Background}
% ─────────────────────────────────────────────────────────────────────────────

\subsection{Shift Register (74HC164)}

A shift register is a chain of D-type flip-flops sharing a common clock line.
On each rising clock edge, the data bit at the serial input is stored in the first
flip-flop, while all other flip-flops take over the value of their predecessor.
The 74HC164 is an 8-bit serial-in / parallel-out shift register.
Its data input is an AND combination of pins A and B; tying both together or holding
one permanently high effectively uses the other as the single serial data input.
An active-low asynchronous clear input (\overline{CLR}) resets all flip-flops to
zero regardless of the clock.

\subsection{D-type Flip-Flop (74HC74)}

The 74HC74 contains two independent positive-edge-triggered D-type flip-flops, each
with active-low asynchronous preset (\overline{PR}) and clear (\overline{CLR}) inputs.
On a rising clock edge, the output $Q$ takes the value present at input $D$.
The asynchronous inputs override the clock: \overline{PR} low forces $Q$ high;
\overline{CLR} low forces $Q$ low, regardless of clock or $D$.

\subsection{Power-On Reset}

After power-up, flip-flops settle in undefined states.
A power-on reset (POR) circuit generates a short active reset pulse immediately
after $V_{CC}$ rises.
An RC circuit charges slowly; the voltage across the capacitor is fed into a
Schmitt-trigger inverter (74HC132).
While $V_C$ is below the upper Schmitt threshold $U_s$, the inverter output remains
high (reset active).
Once $V_C$ exceeds $U_s$, the inverter output falls low (reset released).
The switch-on delay $t_e = \tau \cdot \ln\!\left(\frac{V_{CC}}{V_{CC} - U_s}\right)$
must be longer than the settling time of the supply voltage.
With $R = \SI{1}{\mega\ohm}$ and $C = \SI{1}{\micro\farad}$:
\[
  \tau = R \cdot C = \SI{1}{\mega\ohm} \times \SI{1}{\micro\farad} = \SI{1}{\second}
\]

\subsection{Monoflop (74HC221)}

The 74HC221 is a dual retriggerable monoflop with active-low clear.
It generates a precisely timed output pulse whose duration is set by an external
resistor $R_{ext}$ and capacitor $C_{ext}$:
\begin{equation}
  t_W = 0.7 \cdot R_{ext} \cdot C_{ext}
  \label{eq:monoflop}
\end{equation}

The monoflop is triggered by a rising edge at input~B (with A low) or a falling
edge at input~A (with B high).
During the pulse, the output $\overline{Q}$ is low; after the pulse, it returns high.

\subsection{BCD Decade Counter (74HC160)}

The 74HC160 is a synchronous 4-bit decade counter that counts from 0 to 9 in
natural binary (BCD), then resets to 0 on the next clock edge.
It has an active-low asynchronous clear (\overline{CLR}), a synchronous load input
(\overline{LOAD}), and two count-enable inputs (ENP, ENT).
The ripple-carry output (RCO) goes high when the counter reaches 9, which can be
used to cascade multiple counters.

% ─────────────────────────────────────────────────────────────────────────────
\section{Procedures}
% ─────────────────────────────────────────────────────────────────────────────

All three tasks were carried out according to the laboratory instruction sheet
\textit{Lab\_02\_SequentialCircuits}, Rev.\,1.0, and the design guide
\textit{Lab\_02\_SRC\_Guide}, Rev.\,1.1, both dated 05/06.05.2026.

For each task the following steps were followed:
\begin{enumerate}[leftmargin=1.5em]
  \item Develop and draw the schematic based on the guide.
  \item Simulate the circuit in LTSpice using the provided model files
        (74hc.lib). Note: 74HC160 was replaced by 74HC161 in simulation only
        (identical pinout; counts to 15 instead of 9).
        For 74HC221 the SpiceLine parameter was changed from
        \texttt{TRIPDT=1e-9} to \texttt{TRIPDT=1e-1} to obtain a
        reasonable simulation time.
  \item Present the schematic to the tutor for approval.
  \item Assemble the circuit on the breadboard and verify functionality.
\end{enumerate}

The clock signal was set to \SI{1}{\hertz}, \SI{5}{\volt_{pp}}, \SI{2.5}{\volt} offset.
All ICs were supplied with \SI{5}{\volt} DC.
A 74HC132 Schmitt-trigger inverter was inserted between the waveform generator output
and the clock input of the logic circuit.

% ─────────────────────────────────────────────────────────────────────────────
\section{Results and Discussion}
% ─────────────────────────────────────────────────────────────────────────────

% ════════════════════════════════════════════════════════════════════
\subsection{Task 3.1 -- Chaser Light (8-bit Shift Register)}
% ════════════════════════════════════════════════════════════════════

\subsubsection{Circuit Description}

The chaser light circuit uses the 74HC164 8-bit shift register as its core.
One LED is connected to each of the eight parallel outputs (QA--QH).
A single logic ``1'' is loaded into the first stage at startup and then shifted
one position per clock cycle, illuminating each LED in sequence.

To ensure that exactly one ``1'' is injected at startup and then never again,
one flip-flop of the 74HC74 is used as follows:
\begin{itemize}[leftmargin=1.5em]
  \item The \overline{CLR} input is connected to the power-on reset signal
        (active low during startup, then high).
  \item The D input is held permanently high ($V_{CC}$).
  \item The clock input receives the same clock signal as the shift register.
\end{itemize}
At startup, \overline{CLR} is low, forcing $Q = 0$ regardless of the clock.
At the first rising clock edge after the POR signal releases (goes high), $Q$
takes the value of $D$ = high for exactly one clock period.
Before the second rising edge, $Q$ has returned to low (since \overline{PR} is high
and the flip-flop operates normally).
This single ``1'' pulse at the shift register data input is thus shifted through
all eight stages cyclically.

To ensure the light restarts after LED8, the output of the last flip-flop (QH)
is wired back to the serial data input (A/B) of the 74HC164.

\textbf{Power-on reset:} An RC circuit with $R = \SI{1}{\mega\ohm}$ and
$C = \SI{1}{\micro\farad}$ feeds a 74HC132 Schmitt-trigger inverter.
The switch-on delay is:
\[
  t_e \approx \SI{1}{\second}
\]
which is sufficient to allow the supply voltage to settle before the reset is released.

\subsubsection{LTSpice Simulation}

\begin{figure}[H]
  \centering
  % \includegraphics[width=0.85\linewidth]{fig/3_1_simulation.png}
  \fbox{\rule{0pt}{5.5cm}\rule{12cm}{0pt}}
  \caption{LTSpice simulation -- chaser light: clock, power-on reset, shift register
           data input, and outputs QA--QH.}
  \label{fig:sim_chaser}
\end{figure}

\hint{
  \begin{itemize}[leftmargin=1.2em, itemsep=1pt]
    \item Verify that exactly one output is high at any time
    \item Confirm that after QH goes low, QA goes high on the next clock edge
    \item Check that the power-on reset holds all outputs low until released
    \item Note the gate delay between the POR release and the first ``1'' appearing
          at the data input (one gate propagation delay of the 74HC74)
  \end{itemize}
}

\enter{Describe the simulation result. Confirm the timing: how many clock cycles
until the ``1'' has traversed all 8 stages and wraps around?}

\subsubsection{Breadboard Verification}

\begin{figure}[H]
  \centering
  % \includegraphics[width=0.85\linewidth]{fig/3_1_board.jpg}
  \fbox{\rule{0pt}{5cm}\rule{12cm}{0pt}}
  \caption{Assembled chaser light circuit on the breadboard (photo).}
  \label{fig:board_chaser}
\end{figure}

\enter{Describe the observed behaviour. Did the single LED advance correctly from
LED1 to LED8 and wrap around? Did the power-on reset work (LED1 was the first to
light up)?}

% ════════════════════════════════════════════════════════════════════
\subsection{Task 3.2 -- Catch Circuit (Monoflop Extension)}
% ════════════════════════════════════════════════════════════════════

\subsubsection{Circuit Description}

Task~3.2 extends the chaser light with a ``catch'' function: whenever LED5 is ON and
the push-button is pressed, the moving light freezes at LED5 for \SI{2}{\second},
after which it continues moving normally.

Three sub-circuits are added:

\paragraph{Button Debounce and Catch Detection (74HC74, second flip-flop).}
The push-button is debounced using the second D-type flip-flop of the 74HC74:
\begin{itemize}[leftmargin=1.5em]
  \item \overline{PR}: connected to $V_{CC}$ (not used asynchronously).
  \item \overline{CLR}: connected to the output of LED5's shift register stage
        (Q5 of 74HC164). When LED5 is OFF (Q5 low), \overline{CLR} is low,
        holding $Q_{FF2} = 0$ and preventing any trigger.
  \item D: connected to $V_{CC}$.
  \item Clock (CK): connected to the push-button (via pull-down resistor \SI{10}{\kilo\ohm}).
\end{itemize}
Only when LED5 is ON (Q5 high, \overline{CLR} high) can the flip-flop respond to
a rising edge from the button.
At the first rising edge of the button signal, $Q_{FF2}$ goes high and triggers
the monoflop.
Once the monoflop pulse ends and LED5 goes off (next clock cycle), \overline{CLR}
returns low and resets $Q_{FF2}$, preventing re-triggering.

\paragraph{Monoflop (74HC221).}
The 74HC221 monoflop generates the \SI{2}{\second} hold pulse.
With $R_{ext} = \SI{2}{\mega\ohm}$ and $C_{ext} = \SI{1}{\micro\farad}$:
\begin{equation}
  t_W = 0.7 \times \SI{2}{\mega\ohm} \times \SI{1}{\micro\farad} = \SI{1.4}{\second}
  \label{eq:monoflop_calc}
\end{equation}

\hint{
  \begin{itemize}[leftmargin=1.2em, itemsep=1pt]
    \item The guide specifies $t_W \approx \SI{2}{\second}$; check actual component
          values used and recalculate $t_W$
    \item The monoflop is triggered by $Q_{FF2}$ going high (rising edge at input B,
          with A tied low)
    \item During the pulse, $\overline{Q}$ of the monoflop is low
  \end{itemize}
}

\[
  \boxed{t_{W,\text{calc}} = \SI{1.4}{\second}}
\]

\paragraph{Clock Blocking (74HC00 NAND Gate).}
During the monoflop pulse ($\overline{Q}_{mono}$ low), the clock signal to the shift
register is blocked using a NAND gate from the 74HC00:
\[
  \text{CLK}_{SR} = \overline{\text{CLK}_{FG} \cdot \overline{\overline{Q}_{mono}}}
\]
When $\overline{Q}_{mono}$ is low, one input of the NAND is forced low, its output
is permanently high, and no rising edges reach the shift register clock input.
The light dot therefore stays fixed at LED5 for the duration of the monoflop pulse.

\subsubsection{LTSpice Simulation}

\begin{figure}[H]
  \centering
  % \includegraphics[width=0.85\linewidth]{fig/3_2_simulation.png}
  \fbox{\rule{0pt}{5.5cm}\rule{12cm}{0pt}}
  \caption{LTSpice simulation -- catch circuit: clock, LED5 output, button press,
           monoflop output $\overline{Q}$, and blocked shift register clock.}
  \label{fig:sim_catch}
\end{figure}

\hint{
  \begin{itemize}[leftmargin=1.2em, itemsep=1pt]
    \item Verify that the monoflop only triggers when LED5 is ON
    \item Confirm that the shift register clock is blocked for the duration of $t_W$
    \item Check that pressing the button while LED5 is OFF has no effect
    \item Note: in LTSpice use \texttt{TRIPDT=1e-1} for the 74HC221 SpiceLine
  \end{itemize}
}

\enter{Describe the simulation result. Was the monoflop triggered correctly?
Did the clock blocking work as expected? How long was the measured pause?}

\subsubsection{Breadboard Verification}

\begin{figure}[H]
  \centering
  % \includegraphics[width=0.85\linewidth]{fig/3_2_board.jpg}
  \fbox{\rule{0pt}{5cm}\rule{12cm}{0pt}}
  \caption{Assembled catch circuit on the breadboard (photo).}
  \label{fig:board_catch}
\end{figure}

\enter{Describe observed behaviour. Did LED5 stay on when the button was pressed
during LED5? Did the light resume correctly after the hold time? What happened when
the button was pressed on a different LED?}

% ════════════════════════════════════════════════════════════════════
\subsection{Task 3.3 -- Catch Counter (BCD Counter Extension)}
% ════════════════════════════════════════════════════════════════════

\subsubsection{Circuit Description}

Task~3.3 extends the catch circuit with a BCD decade counter (74HC160) and
4~additional LEDs to count and display the number of successful catches.

Each time the monoflop is triggered (i.e.\ each time the light is successfully
caught), the rising edge of the monoflop's $Q$ output increments the counter by one.
The four parallel BCD outputs (QA--QD) of the 74HC160 drive four LEDs directly,
displaying the catch count in binary (BCD) from 0 to 9.

\textbf{Counter connections:}
\begin{itemize}[leftmargin=1.5em]
  \item \overline{CLR}: connected to the power-on reset signal (active low at startup).
  \item ENP, ENT: both tied to $V_{CC}$ (count always enabled).
  \item \overline{LOAD}: tied to $V_{CC}$ (no parallel load).
  \item CLK: connected to $Q$ of the monoflop (rising edge = one catch).
  \item QA--QD: connected to 4~LEDs via current-limiting resistors.
\end{itemize}

After 9~catches, the counter resets automatically to 0 on the next catch
(decade counter behaviour).

\subsubsection{LTSpice Simulation}

\begin{figure}[H]
  \centering
  % \includegraphics[width=0.85\linewidth]{fig/3_3_simulation.png}
  \fbox{\rule{0pt}{5.5cm}\rule{12cm}{0pt}}
  \caption{LTSpice simulation -- catch counter: monoflop Q output and counter
           outputs QA--QD (BCD code).}
  \label{fig:sim_counter}
\end{figure}

\hint{
  \begin{itemize}[leftmargin=1.2em, itemsep=1pt]
    \item Use 74HC161 in LTSpice instead of 74HC160 (identical pinout, counts to 15)
    \item Verify that QA--QD increment correctly in BCD on each monoflop rising edge
    \item Check that the counter resets to 0 after 9 (hardware) or 15 (simulation)
    \item Confirm that the POR resets the counter to 0 at startup
  \end{itemize}
}

\enter{Describe the simulation result. Did the counter increment correctly?
How was the BCD output verified in simulation?}

\subsubsection{Breadboard Verification}

\begin{figure}[H]
  \centering
  % \includegraphics[width=0.85\linewidth]{fig/3_3_board.jpg}
  \fbox{\rule{0pt}{5cm}\rule{12cm}{0pt}}
  \caption{Assembled catch counter circuit on the breadboard (photo).}
  \label{fig:board_counter}
\end{figure>

\enter{Describe the observed behaviour. Did the 4 LEDs display the correct BCD
count after each catch? Did the counter reset to 0 after 9 catches?
Did power-on reset initialise the counter to 0?}

% ─────────────────────────────────────────────────────────────────────────────
\section{Conclusions}
% ─────────────────────────────────────────────────────────────────────────────

\begin{table}[H]
  \centering
  \caption{Summary of designed and verified circuits.}
  \label{tab:summary}
  \begin{tabular}{llcc}
    \toprule
    Task & Circuit & Simulation & Breadboard \\
    \midrule
    3.1 & Chaser light (74HC164 + 74HC74 + POR) & \enter{OK / NOK} & \enter{OK / NOK} \\
    3.2 & Catch function (+ 74HC221 + 74HC00)   & \enter{OK / NOK} & \enter{OK / NOK} \\
    3.3 & Catch counter (+ 74HC160 + 4 LEDs)    & \enter{OK / NOK} & \enter{OK / NOK} \\
    \bottomrule
  \end{tabular}
\end{table}

\hint{
  \begin{itemize}[leftmargin=1.2em, itemsep=1pt]
    \item Summarise the key findings for all three tasks
    \item Comment on differences between simulation and real hardware (timing,
          propagation delays, power-on behaviour)
    \item Discuss the role of the Schmitt-trigger inverter at the clock input
    \item Discuss the role of the bypass capacitor near IC supply pins
    \item State whether all objectives were met
  \end{itemize}
}

\enter{Write 4--6 sentences summarising the main findings. Refer to
Table~\ref{tab:summary}.}

% ─────────────────────────────────────────────────────────────────────────────
\section{Summary}
% ─────────────────────────────────────────────────────────────────────────────

\enter{3--5 sentence recap: restate the objective, briefly name the methods used,
state the key results, and give the main conclusion. No new information may appear
here.}

% ─────────────────────────────────────────────────────────────────────────────
\section{Outlook}
% ─────────────────────────────────────────────────────────────────────────────

\enter{Suggest 2--3 possible extensions or improvements, for example:
bidirectional chaser light using a multiplexer to select shift direction;
catch counter with 7-segment display instead of binary LEDs;
variable catch time adjustable via potentiometer on the monoflop RC network.}

% ─────────────────────────────────────────────────────────────────────────────
\begin{thebibliography}{9}
% ─────────────────────────────────────────────────────────────────────────────

\bibitem{Floyd2006}
T.\,Floyd,
\textit{Digital Fundamentals}, 9th~ed.
Pearson Prentice Hall, 2006, ch.~7.

\bibitem{TietzeSchenk1999}
U.\,Tietze and C.\,Schenk,
\textit{Halbleiterschaltungstechnik}, 11.~Auflage.
Springer, 1999.

\bibitem{ReportGuide}
C.\,Thorrold,
\textit{Guide for Writing Lab Reports}, Version~1.0.

\bibitem{74HC164_DS}
Texas Instruments,
\textit{74HC164 8-Bit Serial-In, Parallel-Out Shift Register -- Datasheet},
\url{https://www.ti.com/lit/ds/symlink/sn74hc164.pdf}, accessed May 2026.

\bibitem{74HC74_DS}
Texas Instruments,
\textit{74HC74 Dual D-Type Positive-Edge-Triggered Flip-Flop -- Datasheet},
\url{https://www.ti.com/lit/ds/symlink/sn74hc74.pdf}, accessed May 2026.

\bibitem{74HC221_DS}
Texas Instruments,
\textit{74HC221 Dual Monostable Multivibrator -- Datasheet},
\url{https://www.ti.com/lit/ds/symlink/cd74hc221.pdf}, accessed May 2026.

\bibitem{74HC160_DS}
Texas Instruments,
\textit{74HC160 Synchronous 4-Bit Decade Counter -- Datasheet},
\url{https://www.ti.com/lit/ds/symlink/sn74hc160.pdf}, accessed May 2026.

\end{thebibliography}

% ─────────────────────────────────────────────────────────────────────────────
\appendix
\section{Appendix A -- Schematics}
% ─────────────────────────────────────────────────────────────────────────────

\subsection{Task 3.1 -- Chaser Light Schematic}

\begin{figure}[H]
  \centering
  % \includegraphics[width=0.9\linewidth]{fig/schematic_3_1.png}
  \fbox{\rule{0pt}{7cm}\rule{14cm}{0pt}}
  \caption{Complete schematic -- chaser light circuit (Task 3.1).}
  \label{fig:schematic_31}
\end{figure}

\subsection{Task 3.2 -- Catch Circuit Schematic}

\begin{figure}[H]
  \centering
  % \includegraphics[width=0.9\linewidth]{fig/schematic_3_2.png}
  \fbox{\rule{0pt}{7cm}\rule{14cm}{0pt}}
  \caption{Complete schematic -- catch circuit (Task 3.2).}
  \label{fig:schematic_32}
\end{figure}

\subsection{Task 3.3 -- Catch Counter Schematic}

\begin{figure}[H]
  \centering
  % \includegraphics[width=0.9\linewidth]{fig/schematic_3_3.png}
  \fbox{\rule{0pt}{7cm}\rule{14cm}{0pt}}
  \caption{Complete schematic -- catch counter circuit (Task 3.3).}
  \label{fig:schematic_33}
\end{figure}

\section{Appendix B -- LTSpice Simulation Screenshots}
% ─────────────────────────────────────────────────────────────────────────────

\enter{Insert full LTSpice waveform screenshots here, one per task, with all
relevant signals labelled. Note which node names were probed.}

\end{document}